


[[1.1]]
Exercise 1: Dirac and Klein-Gordon equation(4 points)
Let $\psi$ be a Dirac spinor solving the Dirac equation
$$
\left[\mathrm{i} \gamma^\mu \partial_\mu-m\right] \psi(x)=0
$$
Show that it is also a solution of the Klein-Gordon equation
$$
\left[\partial_\mu \partial^\mu+m^2\right] \psi(x)=0
$$



Exercise 2: Bilinear covariants $\mathbf{I}(3+3=6$ points $)$

[[2.1]]
i) First show that
$$
\left\{\gamma^\mu, \gamma^\nu\right\}=2 g^{\mu \nu}
$$
and use your results to show that
$$
\gamma^{\mu \dagger}=\gamma^0 \gamma^\mu \gamma^0
$$

[[2.2]]
ii) In the lecture you found the following Lorentz transformation of a Dirac spinor $\psi$
$$
\psi(x) \rightarrow \psi^{\prime}\left(x^{\prime}\right)=S(\Lambda) \psi(x), \quad S(\Lambda)=e^{-\frac{i}{4} \omega^{\mu \nu} \sigma_{\mu \nu}}
$$
with $\sigma_{\mu \nu}=\frac{i}{2}\left[\gamma_\mu, \gamma_\nu\right]$.
Furthermore you have shown that $\bar{\psi} \psi$ is a Lorentz scalar, therefore the inverse transformation is given by
$$
\bar{\psi}(x) \rightarrow \bar{\psi}^{\prime}\left(x^{\prime}\right) S^{-1}(\Lambda), \quad S^{-1}(\Lambda)=e^{\frac{i}{4} \omega^{\mu \nu} \sigma_{\mu \nu}}
$$
Show that $S^{-1}(\Lambda)$ can also be written as
$$
S^{-1}(\Lambda)=\left(\gamma^0 S(\Lambda) \gamma^0\right)^{\dagger}
$$
(instead of $S^{-1}(\Lambda)=S^{\dagger}(\Lambda)$ as one would expect naively).
Hint: Use the result of $i$ )



[[3]]
Recall the infinitesimal representation of the Lorentz transformation
$$
\Lambda_\nu^\mu=\delta_\nu^\mu+\delta \omega_\nu^\mu
$$
and the Lorentz transformation of a spinor
$$
S(\Lambda)=1-\frac{\mathrm{i}}{4} \sigma_{\mu \nu} \delta \omega^{\mu \nu}+o\left(\delta \omega^2\right), \quad \sigma_{\mu \nu}=\frac{\mathrm{i}}{2}\left[\gamma_\mu, \gamma_\nu\right]
$$

[[3.1]]
i) Using these transformations, show that the equation
$$
S^{-1}(\Lambda) \gamma^\mu S(\Lambda)=\Lambda_\nu^\mu \gamma^\nu
$$
familiar from the lecture, is satisfied up to first order in $\delta \omega$.

[[3.2]]
ii) In the lecture you showed that $\bar{\psi} \psi$ transforms as a Lorentz scalar and $\bar{\psi} \gamma^\mu \psi$ transforms as a Lorentz vector. Furthermore you introduced an additional gamma matrix $\gamma^5=\mathrm{i} \gamma^0 \gamma^1 \gamma^2 \gamma^3$ with the transformation property
$$
S^{-1}(\Lambda) \gamma^5 S(\Lambda)=\operatorname{det}(\Lambda) \gamma^5
$$
Prove the behaviour under Lorentz transformation of the following bilinear covariants
$$
\begin{aligned}
\bar{\psi} \gamma^5 \psi & \rightarrow \operatorname{det}(\Lambda) \bar{\psi} \gamma^5 \psi & & \text { pseudoscalar } \\
\bar{\psi} \gamma^\mu \gamma^5 \psi & \rightarrow \operatorname{det}(\Lambda) \Lambda_\nu^\mu \bar{\psi} \gamma^\nu \gamma^5 \psi & & \text { axial vector } \\
\bar{\psi} \sigma^{\mu \nu} \psi & \rightarrow \Lambda_\alpha^\mu \Lambda_\beta^\nu \bar{\psi} \sigma^{\alpha \beta} \psi & & \text { antisymmetric tensor }
\end{aligned}
$$
Hint: Equation (10) will be very useful.


[[4]]
Exercise 4: Continuity equation in relativistic quantum mechanics II(4 points) Let the Dirac spinor $\psi$ be a solution of the Dirac equation
$$
\left[\mathrm{i} \gamma^\mu \partial_\mu-m\right] \psi(x)=0
$$

[[4.1]]
Show that the current
$$
j^\mu(x)=\bar{\psi} \gamma^\mu \psi
$$
satisfies a continuity equation
$$
\partial_\mu j^\mu(x)=0
$$

[[4.2]]
Discuss wether $j^0=\rho$ is positive definite in the case of the Dirac equation.